\begin{longtable} {>{\raggedright\arraybackslash}p{0.3\linewidth}
    >{\raggedright\arraybackslash}p{0.7\linewidth}}
    \caption{Übersicht untersuchter Anwendungsszenarien im MIR} \label{tab:scenario} \\

    \toprule
    \textbf{Kategorie} & \textbf{Beschreibung} \\
    \midrule
    \endfirsthead

    \toprule
    \textbf{Kategorie} & \textbf{Beschreibung} \\
    \midrule
    \endhead

    \midrule
    \multicolumn{2}{r}{\textit{Fortsetzung auf nächster Seite}} \\
    \endfoot

    \bottomrule
    \endlastfoot

% --- EINTRAG ---
    \textbf{Szenario} & Waldbaden \\
    \textbf{Inhaltstyp} & vorgerendert \\
    \textbf{Verwendete Technologie} & 360° Sound, 270° Visuals \\
    \textbf{Szenarioziel} & atmosphärisch (ambient) \\
    \textbf{Interaktionsform} & keine \\
    \textbf{Systemreaktion} & keine \\
    \textbf{Interaktionstyp} & keine \\
    \textbf{Navigation} & fester Perspektivwechsel \\
    \textbf{Physische Inszenierungselemente} & Wiesenboden \\
    \textbf{Beschreibung Interaktion} & Keine Reaktion vom System auf Nutzer:innen; Nutzer:innen sollten sich möglichst nicht hinstellen/laufen, weil Bodenerhöhung Immersion im oberen Sichtfeld beeinflusst \\[0.8em]
    \midrule

    % --- EINTRAG ---
    \textbf{Szenario} & Fußballstadion \\
    \textbf{Inhaltstyp} & vorgerendert \\
    \textbf{Verwendete Technologie} & 360° Sound, 270° Visuals, Bodenvibration \\
    \textbf{Szenarioziel} & atmosphärisch (narrativ) \\
    \textbf{Interaktionsform} & keine \\
    \textbf{Systemreaktion} & keine \\
    \textbf{Interaktionstyp} & keine \\
    \textbf{Navigation} & fester Perspektivwechsel \\
    \textbf{Physische Inszenierungselemente} & Stühle \\
    \textbf{Beschreibung Interaktion} & Keine Reaktion vom System auf Nutzer:innen \\[0.8em]
    \midrule

    % --- EINTRAG ---
    \textbf{Szenario} & Kokon \\
    \textbf{Inhaltstyp} & vorgerendert \\
    \textbf{Verwendete Technologie} & 360° Sound, 270° Visuals, Bodenvibration \\
    \textbf{Szenarioziel} & atmosphärisch (ambient) \\
    \textbf{Interaktionsform} & keine \\
    \textbf{Systemreaktion} & keine \\
    \textbf{Interaktionstyp} & keine \\
    \textbf{Navigation} & keine \\
    \textbf{Physische Inszenierungselemente} & Tücher \\
    \textbf{Beschreibung Interaktion} & Keine Reaktion vom System auf Nutzer:innen; Nutzer:innen sollten sich möglichst nicht hinstellen/laufen, weil Bodenerhöhung Immersion im oberen Sichtfeld beeinflusst \\[0.8em]
    \midrule

    % --- EINTRAG ---
    \textbf{Szenario} & Zug fahren \\
    \textbf{Inhaltstyp} & vorgerendert \\
    \textbf{Verwendete Technologie} & 270° Visuals \\
    \textbf{Szenarioziel} & atmosphärisch (ambient) \\
    \textbf{Interaktionsform} & keine \\
    \textbf{Systemreaktion} & zeitbasierter Wechsel der Szene (beispielsweise 25 Minuten Zug fahren, dann Szenenwechsel zu 5 Minuten in der Natur) \\
    \textbf{Interaktionstyp} & keine \\
    \textbf{Navigation} & fester Perspektivwechsel \\
    \textbf{Physische Inszenierungselemente} & Stüfle \\
    \textbf{Beschreibung Interaktion} & Keine Reaktion vom System auf Nutzer:innen; Nach bestimmter Zeit wechselt Szene, sodass Nutzer:innen aufstehen und sich bewegen sollen (Arbeitspause vom Zug fahren) \\[0.8em]
    \midrule

    % --- EINTRAG ---
    \textbf{Szenario} & YouTube 360° \\
    \textbf{Inhaltstyp} & vorgerendert \\
    \textbf{Verwendete Technologie} & 360° Sound, 270° Visuals \\
    \textbf{Szenarioziel} & atmosphärisch (narrativ) \\
    \textbf{Interaktionsform} & keine \\
    \textbf{Systemreaktion} & keine \\
    \textbf{Interaktionstyp} & keine \\
    \textbf{Navigation} & fester Perspektivwechsel \\
    \textbf{Physische Inszenierungselemente} & Stühle \\
    \textbf{Beschreibung Interaktion} & Keine Reaktion vom System auf Nutzer:innen \\[0.8em]
    \midrule

    % --- EINTRAG ---
    \textbf{Szenario} & Jungle Level \\
    \textbf{Inhaltstyp} & live gerendert \\
    \textbf{Verwendete Technologie} & 360° Sound, 270° Visuals \\
    \textbf{Szenarioziel} & atmosphärisch (ambient) \\
    \textbf{Interaktionsform} & Tastatur/Maus \\
    \textbf{Systemreaktion} & freie Bewegung im Raum \\
    \textbf{Interaktionstyp} & explizit \\
    \textbf{Navigation} & frei (mit Tastatur) \\
    \textbf{Physische Inszenierungselemente} & keine \\
    \textbf{Interaktionshardware} & TEXT \\
    \textbf{Beschreibung Interaktion} & Freie Bewegung im Jungle Level mit Tastatur; Tastatur ist aber kein Teil des MIR Aufbau, daher eingeschränkt nutzbar \\[0.8em]
    \midrule

    % --- EINTRAG ---
    \textbf{Szenario} & Smart Factory \\
    \textbf{Inhaltstyp} & live gerendert \\
    \textbf{Verwendete Technologie} & 360° Sound, 270° Visuals \\
    \textbf{Szenarioziel} & Simulation \\
    \textbf{Interaktionsform} & Tastatur/Maus, Gesteninteraktion \\
    \textbf{Systemreaktion} & Zustandsänderung im Reparaturarbeitsplatz ändert Systemzustände (Beispielsweise korrektes Ausführen von Reparaturanweisungen stoppt Alarm)\\
    \textbf{Interaktionstyp} & explizit \\
    \textbf{Navigation} & eingeschränkt (mit Tastatur) \\
    \textbf{Physische Inszenierungselemente} & Manueller Reparaturarbeitsplatz \\
    \textbf{Beschreibung Interaktion} & Administrator:innen können vom Rechnera aus ansteuern, dass Szenenelemente Aktionen ausführen, zum Beispiel dass eine Dohne losfliegen und umkehren oder dass Roboter los fahren. Nutzer:innen können am Reparaturarbeitsplatz mithilfe von Gestensteuerung Arbeitsschritte als beendet markieren \\[0.8em]
    \midrule

    % --- EINTRAG ---
    \textbf{Szenario} & OP \\
    \textbf{Inhaltstyp} & vorgerendert \\
    \textbf{Verwendete Technologie} & 360° Sound, 270° Visuals \\
    \textbf{Szenarioziel} & Lehre \\
    \textbf{Interaktionsform} & keine \\
    \textbf{Systemreaktion} & keine \\
    \textbf{Interaktionstyp} & keine \\
    \textbf{Navigation} & fester Perspektivwechsel \\
    \textbf{Physische Inszenierungselemente} & Lehrobjekte \\
    \textbf{Beschreibung Interaktion} & Keine Reaktion vom System auf Nutzer:innen \\[0.8em]
    \midrule

    % --- EINTRAG ---
    \textbf{Szenario} & Hubschrauberlandeplatz \\
    \textbf{Inhaltstyp} & vorgerendert \\
    \textbf{Verwendete Technologie} & 360° Sound, 270° Visuals \\
    \textbf{Szenarioziel} & Lehre \\
    \textbf{Interaktionsform} & keine \\
    \textbf{Systemreaktion} & keine \\
    \textbf{Interaktionstyp} & keine \\
    \textbf{Navigation} & keine \\
    \textbf{Physische Inszenierungselemente} & Lehrobjekte \\
    \textbf{Interaktionshardware} & Keine \\
    \textbf{Beschreibung Interaktion} & Keine Reaktion vom System auf Nutzer:innen \\[0.8em]
    \midrule

\end{longtable}